\section{Implementation Plan}

\subsection{Phase 1: Environment Setup (Day 1)}

\begin{enumerate}[nosep]
    \item Create a Python virtual environment
    \item Install dependencies:
\begin{minted}{bash}
pip install torch transformers datasets evaluate scikit-learn matplotlib
\end{minted}
    \item Verify GPU access:
\begin{minted}{python}
python -c "import torch; print(torch.backends.mps.is_available())"
\end{minted}
\end{enumerate}

\subsection{Phase 2: Dataset Exploration (Day 2)}

\begin{enumerate}[nosep]
    \item Load XNLI for Hindi:
\begin{minted}{python}
from datasets import load_dataset
dataset = load_dataset("xnli", "hi")
print(dataset["train"][0])
\end{minted}
    \item Understand label mapping: 0 = Entailment, 1 = Neutral, 2 = Contradiction
    \item Statistics: 392k training pairs (machine-translated from English), 2.5k validation, 5k test
\end{enumerate}

\subsection{Phase 3: Model \& Tokenizer Setup (Day 3)}

\begin{minted}{python}
from transformers import AutoTokenizer, AutoModelForSequenceClassification

model_name = "xlm-roberta-base"
tokenizer = AutoTokenizer.from_pretrained(model_name)
model = AutoModelForSequenceClassification.from_pretrained(
    model_name, num_labels=3
)
\end{minted}

\subsection{Phase 4: Training Pipeline (Day 4-5)}

\begin{enumerate}[nosep]
    \item Tokenize the dataset (padding, truncation to 128 tokens)
    \item Set device to Metal (MPS) for Apple Silicon:
\begin{minted}{python}
import torch
device = torch.device("mps" if torch.backends.mps.is_available() else "cpu")
model.to(device)
\end{minted}
    \item Configure TrainingArguments:
\begin{minted}{python}
TrainingArguments(
    output_dir="./xnli-hindi-model",
    evaluation_strategy="epoch",
    save_strategy="epoch",
    num_train_epochs=3,
    per_device_train_batch_size=16,
    per_device_eval_batch_size=64,
    warmup_ratio=0.1,
    logging_dir="./logs",
    fp16=True,
    load_best_model_at_end=True,
    metric_for_best_model="accuracy",
)
\end{minted}
    \item Define compute\_metrics function (accuracy, precision, recall, F1)
    \item Run Trainer:
\begin{minted}{python}
trainer = Trainer(
    model=model,
    args=training_args,
    train_dataset=encoded["train"],
    eval_dataset=encoded["validation"],
    tokenizer=tokenizer,
    compute_metrics=compute_metrics,
)
trainer.train()
\end{minted}
    \item Expected training time: 30-45 minutes on M5 Apple GPU (per epoch)
\end{enumerate}

\subsection{Phase 5: Evaluation (Day 6)}

\begin{minted}{python}
results = trainer.evaluate(encoded["test"])
print("Test Accuracy:", results["eval_accuracy"])
\end{minted}

Expected results:
\begin{itemize}[nosep]
    \item \textbf{Accuracy:} 85-88\%
    \item \textbf{F1 Score:} 0.84-0.87
    \item \textbf{Precision:} 0.83-0.87
    \item \textbf{Recall:} 0.85-0.88
\end{itemize}

\subsection{Phase 6: Inference Demo (Day 7)}

Build a reusable inference function:
\begin{minted}{python}
device = torch.device("mps" if torch.backends.mps.is_available() else "cpu")
model.to(device)

def predict_entailment(premise, hypothesis):
    inputs = tokenizer(premise, hypothesis,
                       return_tensors="pt", truncation=True)
    inputs = {k: v.to(device) for k, v in inputs.items()}
    outputs = model(**inputs)
    probs = torch.nn.functional.softmax(outputs.logits, dim=-1)
    label = torch.argmax(probs).item()
    return {0: "Entailment", 1: "Neutral", 2: "Contradiction"}[label]
\end{minted}
